
\documentclass[a4paper,11pt]{article}
\usepackage[a4paper, margin=8em]{geometry}

% usa i pacchetti per la scrittura in italiano
\usepackage[french,italian]{babel}
\usepackage[T1]{fontenc}
\usepackage[utf8]{inputenc}
\frenchspacing 

% usa i pacchetti per la formattazione matematica
\usepackage{amsmath, amssymb, amsthm, amsfonts}

% usa altri pacchetti
\usepackage{gensymb}
\usepackage{hyperref}
\usepackage{standalone}

\usepackage{colortbl}

\usepackage{xstring}
\usepackage{karnaugh-map}

% imposta il titolo
\title{Embedded Operating Systems and Architectures notes}
\author{Luca Seggiani}
\date{2026}

% imposta lo stile
% usa helvetica
\usepackage[scaled]{helvet}
% usa palatino
\usepackage{palatino}
% usa un font monospazio guardabile
\usepackage{lmodern}

\renewcommand{\rmdefault}{ppl}
\renewcommand{\sfdefault}{phv}
\renewcommand{\ttdefault}{lmtt}

% circuiti
\usepackage{circuitikz}
\usetikzlibrary{babel}

% testo cerchiato
\newcommand*\circled[1]{\tikz[baseline=(char.base)]{
            \node[shape=circle,draw,inner sep=2pt] (char) {#1};}}

% disponi il titolo
\makeatletter
\renewcommand{\maketitle} {
	\begin{center} 
		\begin{minipage}[t]{.8\textwidth}
			\textsf{\huge\bfseries \@title} 
		\end{minipage}%
		\begin{minipage}[t]{.2\textwidth}
			\raggedleft \vspace{-1.65em}
			\textsf{\small \@author} \vfill
			\textsf{\small \@date}
		\end{minipage}
		\par
	\end{center}

	\thispagestyle{empty}
	\pagestyle{fancy}
}
\makeatother

% disponi teoremi
\usepackage{tcolorbox}
\newtcolorbox[auto counter, number within=section]{theorem}[2][]{%
	colback=blue!10, 
	colframe=blue!40!black, 
	sharp corners=northwest,
	fonttitle=\sffamily\bfseries, 
	title=Theorem~\thetcbcounter: #2, 
	#1
}

% disponi definizioni
\newtcolorbox[auto counter, number within=section]{definition}[2][]{%
	colback=red!10,
	colframe=red!40!black,
	sharp corners=northwest,
	fonttitle=\sffamily\bfseries,
	title=Definition~\thetcbcounter: #2,
	#1
}

% disponi codice
\usepackage{listings}
\usepackage[table]{xcolor}

\definecolor{codegreen}{rgb}{0,0.6,0}
\definecolor{codegray}{rgb}{0.5,0.5,0.5}
\definecolor{codepurple}{rgb}{0.58,0,0.82}
\definecolor{backcolour}{rgb}{0.95,0.95,0.92}

\lstdefinestyle{codestyle}{
		backgroundcolor=\color{black!5}, 
		commentstyle=\color{codegreen},
		keywordstyle=\bfseries\color{magenta},
		numberstyle=\sffamily\tiny\color{black!60},
		stringstyle=\color{green!50!black},
		basicstyle=\ttfamily\footnotesize,
		breakatwhitespace=false,         
		breaklines=true,                 
		captionpos=b,                    
		keepspaces=true,                 
		numbers=left,                    
		numbersep=5pt,                  
		showspaces=false,                
		showstringspaces=false,
		showtabs=false,                  
		tabsize=2
}

\lstdefinestyle{shellstyle}{
		backgroundcolor=\color{black!5}, 
		basicstyle=\ttfamily\footnotesize\color{black}, 
		commentstyle=\color{black}, 
		keywordstyle=\color{black},
		numberstyle=\color{black!5},
		stringstyle=\color{black}, 
		showspaces=false,
		showstringspaces=false, 
		showtabs=false, 
		tabsize=2, 
		numbers=none, 
		breaklines=true
}

\lstdefinelanguage{x86}{ 
  keywords={AAA, AAD, AAM, AAS, ADC, ADCB, ADCW, ADCL, ADD, ADDB, ADDW, ADDL, AND, ANDB, ANDW, ANDL,
        ARPL, BOUND, BSF, BSFL, BSFW, BSR, BSRL, BSRW, BSWAP, BT, BTC, BTCB, BTCW, BTCL, BTR, 
        BTRB, BTRW, BTRL, BTS, BTSB, BTSW, BTSL, CALL, CBW, CDQ, CLC, CLD, CLI, CLTS, CMC, CMP,
        CMPB, CMPW, CMPL, CMPS, CMPSB, CMPSD, CMPSW, CMPXCHG, CMPXCHGB, CMPXCHGW, CMPXCHGL,
        CMPXCHG8B, CPUID, CWDE, DAA, DAS, DEC, DECB, DECW, DECL, DIV, DIVB, DIVW, DIVL, ENTER,
        HLT, IDIV, IDIVB, IDIVW, IDIVL, IMUL, IMULB, IMULW, IMULL, IN, INB, INW, INL, INC, INCB,
        INCW, INCL, INS, INSB, INSD, INSW, INT, INT3, INTO, INVD, INVLPG, IRET, IRETD, JA, JAE,
        JB, JBE, JC, JCXZ, JE, JECXZ, JG, JGE, JL, JLE, JMP, JNA, JNAE, JNB, JNBE, JNC, JNE, JNG,
        JNGE, JNL, JNLE, JNO, JNP, JNS, JNZ, JO, JP, JPE, JPO, JS, JZ, LAHF, LAR, LCALL, LDS,
        LEA, LEAVE, LES, LFS, LGDT, LGS, LIDT, LMSW, LOCK, LODSB, LODSD, LODSW, LOOP, LOOPE,
        LOOPNE, LSL, LSS, LTR, MOV, MOVB, MOVW, MOVL, MOVSB, MOVSD, MOVSW, MOVSX, MOVSXB,
        MOVSXW, MOVSXL, MOVZX, MOVZXB, MOVZXW, MOVZXL, MUL, MULB, MULW, MULL, NEG, NEGB, NEGW,
        NEGL, NOP, NOT, NOTB, NOTW, NOTL, OR, ORB, ORW, ORL, OUT, OUTB, OUTW, OUTL, OUTSB, OUTSD,
        OUTSW, POP, POPL, POPW, POPB, POPA, POPAD, POPF, POPFD, PUSH, PUSHL, PUSHW, PUSHB, PUSHA, 
        RCLL, RCR, RCRB, RCRW, RCRL, RDMSR, RDPMC, RDTSC, REP, REPE, REPNE, RET, ROL, ROLB, ROLW,
        ROLL, ROR, RORB, RORW, RORL, SAHF, SAL, SALB, SALW, SALL, SAR, SARB, SARW, SARL, SBB,
        SBBB, SBBW, SBBL, SCASB, SCASD, SCASW, SETA, SETAE, SETB, SETBE, SETC, SETE, SETG, SETGE,
        SETL, SETLE, SETNA, SETNAE, SETNB, SETNBE, SETNC, SETNE, SETNG, SETNGE, SETNL, SETNLE,
        SETNO, SETNP, SETNS, SETNZ, SETO, SETP, SETPE, SETPO, SETS, SETZ, SGDT, SHL, SHLB, SHLW,
        SHLL, SHLD, SHR, SHRB, SHRW, SHRL, SHRD, SIDT, SLDT, SMSW, STC, STD, STI, STOSB, STOSD,
        STOSW, STR, SUB, SUBB, SUBW, SUBL, TEST, TESTB, TESTW, TESTL, VERR, VERW, WAIT, WBINVD,
        XADD, XADDB, XADDW, XADDL, XCHG, XCHGB, XCHGW, XCHGL, XLAT, XLATB, XOR, XORB, XORW, XORL},
  keywordstyle=\color{blue}\bfseries,
  ndkeywordstyle=\color{darkgray}\bfseries,
  identifierstyle=\color{black},
  sensitive=false,
  comment=[l]{\#},
  morecomment=[s]{/*}{*/},
  commentstyle=\color{purple}\ttfamily,
  stringstyle=\color{red}\ttfamily,
  morestring=[b]',
  morestring=[b]"
}

\lstdefinelanguage{riscv}{
  keywords={
    LUI, AUIPC,
    JAL, JALR,
    BEQ, BNE, BLT, BGE, BLTU, BGEU,
    LB, LH, LW, LBU, LHU, LD,
    SB, SH, SW, SD,
    ADDI, SLTI, SLTIU, XORI, ORI, ANDI,
    SLLI, SRLI, SRAI,
    ADD, SUB, SLL, SLT, SLTU, XOR, SRL, SRA, OR, AND,
    ADDIW, SLLIW, SRLIW, SRAIW,
    ADDW, SUBW, SLLW, SRLW, SRAW,
    MUL, MULH, MULHSU, MULHU, DIV, DIVU, REM, REMU,
    MULW, DIVW, DIVUW, REMW, REMUW,
    LR.W, SC.W, AMOSWAP.W, AMOADD.W, AMOXOR.W,
    AMOAND.W, AMOOR.W, AMOMIN.W, AMOMAX.W,
    AMOMINU.W, AMOMAXU.W,
    LR.D, SC.D, AMOSWAP.D, AMOADD.D, AMOXOR.D,
    AMOAND.D, AMOOR.D, AMOMIN.D, AMOMAX.D,
    AMOMINU.D, AMOMAXU.D,
    FLW, FSW,
    FMADD.S, FMSUB.S, FNMSUB.S, FNMADD.S,
    FADD.S, FSUB.S, FMUL.S, FDIV.S, FSQRT.S,
    FSGNJ.S, FSGNJN.S, FSGNJX.S,
    FMIN.S, FMAX.S,
    FCVT.W.S, FCVT.WU.S, FCVT.S.W, FCVT.S.WU,
    FMV.X.W, FMV.W.X,
    FEQ.S, FLT.S, FLE.S,
    FLD, FSD,
    FMADD.D, FMSUB.D, FNMSUB.D, FNMADD.D,
    FADD.D, FSUB.D, FMUL.D, FDIV.D, FSQRT.D,
    FSGNJ.D, FSGNJN.D, FSGNJX.D,
    FMIN.D, FMAX.D,
    FCVT.W.D, FCVT.WU.D, FCVT.D.W, FCVT.D.WU,
    FCVT.S.D, FCVT.D.S,
    FMV.X.D, FMV.D.X,
    FEQ.D, FLT.D, FLE.D,
    CSRRW, CSRRS, CSRRC,
    CSRRWI, CSRRSI, CSRRCI,
    FENCE.I,
    FENCE,
    ECALL, EBREAK,
    MRET, SRET, URET,
    WFI,
    SFENCE.VMA,
    NOP, LI, LA, MV,
    NOT, NEG, NEGW,
    SEQZ, SNEZ, SLTZ, SGTZ,
    BEQZ, BNEZ, BLEZ, BGEZ, BLTZ, BGTZ,
    BGT, BLE, BGTU, BLEU,
    J, JR, RET,
    CALL, TAIL,
    RDINSTRET, RDCYCLE, RDTIME,
    CSRR, CSRW, CSRS, CSRC,
    CSRWI, CSRSI, CSRCI
  },
  keywordstyle=\color{blue}\bfseries,
  ndkeywordstyle=\color{darkgray}\bfseries,
  identifierstyle=\color{black},
  sensitive=false,
  comment=[l]{\#},
  morecomment=[s]{/*}{*/},
  commentstyle=\color{purple}\ttfamily,
  stringstyle=\color{red}\ttfamily,
  morestring=[b]',
  morestring=[b]"
}

\lstdefinelanguage{arm}{
  keywords={
    ADC, ADD, AND, BIC, CMN, CMP, EOR, MOV, MVN,
    RSB, RSC, SBC, SUB, TST, TEQ,
    MLA, MLS, MUL, SMLAL, SMULL, UMLAL, UMULL,
    QADD, QSUB, QDADD, QDSUB,
    QADD16, QADD8, QASX, QSAX,
    QSUB16, QSUB8, QSAX, QSUBADDX,
    ASR, LSL, LSR, ROR, RRX,
    B, BL, BX, BLX,
    BXJ,
    BEQ, BNE, BCS, BHS, BCC, BLO,
    BMI, BPL, BVS, BVC,
    BHI, BLS, BGE, BLT, BGT, BLE,
    BAL,
    LDR, LDRB, LDRH, LDRSB, LDRSH,
    LDRD, LDM, LDMIA, LDMIB, LDMDA, LDMDB,
    LDRT, LDRBT, LDRHT, LDRSBT, LDRSHT,
    STR, STRB, STRH, STRD,
    STM, STMIA, STMIB, STMDA, STMDB,
    STRT, STRBT, STRHT,
    PUSH, POP,
    SWP, SWPB,
    REV, REV16, REVSH,
    BFC, BFI, BFXIL, SBFX, UBFX,
    CLZ,
    LDREX, LDREXB, LDREXH, LDREXD,
    STREX, STREXB, STREXH, STREXD,
    DMB, DSB, ISB,
    MCR, MCRR, MRC, MRRC,
    MRS, MSR,
    CPS, SETEND, SVC, BKPT,
    CLREX, WFE, WFI, SEV,
    YIELD, NOP,
    VLDR, VSTR,
    VMOV, VABS, VNEG,
    VADD, VSUB, VMUL, VDIV,
    VMLA, VMLS, VNMLA, VNMLS, VNMUL,
    VMAX, VMIN,
    VSQRT,
    VCMP, VCMPE,
    VCVT,
    VCLZ, VCLS,
    VAND, VORR, VEOR, VBIC, VMVN,
    VLD1, VLD2, VLD3, VLD4,
    VST1, VST2, VST3, VST4,
    ADR, ADCS, ADDS, SUBS, MOVS, ANDS, ORRS, EORS
  },
  keywordstyle=\color{blue}\bfseries,
  ndkeywordstyle=\color{darkgray}\bfseries,
  identifierstyle=\color{black},
  sensitive=false,
  comment=[l]{@},
  morecomment=[l]{;},
  morecomment=[s]{/*}{*/},
  commentstyle=\color{purple}\ttfamily,
  stringstyle=\color{red}\ttfamily,
  morestring=[b]',
  morestring=[b]"
}

\lstset{style=codestyle, language=C}

% disponi sezioni
\usepackage{titlesec}

\titleformat{\section}
	{\sffamily\Large\bfseries} 
	{\thesection}{1em}{} 
\titleformat{\subsection}
	{\sffamily\large\bfseries}   
	{\thesubsection}{1em}{} 
\titleformat{\subsubsection}
	{\sffamily\normalsize\bfseries} 
	{\thesubsubsection}{1em}{}

% tikz
\usepackage{tikz}

% float
\usepackage{float}

% grafici
\usepackage{pgfplots}
\pgfplotsset{width=10cm,compat=1.9}

% disponi alberi
\usepackage{forest}

\forestset{
	rectstyle/.style={
		for tree={rectangle,draw,font=\large\sffamily}
	},
	roundstyle/.style={
		for tree={circle,draw,font=\large}
	}
}

% disponi algoritmi
\usepackage{algorithm}
\usepackage{algorithmic}
\makeatletter
\renewcommand{\ALG@name}{Algorithm}
\makeatother

% disponi numeri di pagina
\usepackage{fancyhdr}
\fancyhf{} 
\fancyfoot[L]{\sffamily{\thepage}}

\makeatletter
\fancyhead[L]{\raisebox{1ex}[0pt][0pt]{\sffamily{\@title \ \@date}}} 
\fancyhead[R]{\raisebox{1ex}[0pt][0pt]{\sffamily{\@author}}}
\makeatother

\begin{document}
% sezione (data)
\section{Lesson 30-09-26}

% stili pagina
\thispagestyle{empty}
\pagestyle{fancy}

% testo

\subsection{Memory}

Last week we went over the structure of an embedded system MCU.
Now we focus on memory, with the target of defining a memory map and defining a linker script for it.

\subsubsection{Types of memory}

Let's quickly go over the \textit{types} of memory available to us.

\par\smallskip
\textbf{\textsf{SRAM}}

\begin{itemize}

\item
  \textbf{SRAM} (static RAM) sits on the CPU's memory bus: every byte
  has an address, and the CPU reads or writes it with a normal
  load/store;
\item
  \textbf{Fast}: on a microcontroller, one access takes about one clock
  cycle;
\item
  \textbf{Byte-writable}: we can change a single byte, any time, as
  often as we like;
\item
  \textbf{Volatile}: power off and everything is lost; at power-on the
  content is \textbf{undefined} (random bits, not zeros).
\end{itemize}

SRAM is where anything that \textbf{changes} while the program runs must
live: variables, the stack.

\par\smallskip
\textbf{\textsf{NOR flash}}

\begin{itemize}

\item
  Also sits on the memory bus, like a ROM: every word has an address,
  random-access reads are fast;
\item
	So the CPU can \textbf{fetch} instructions directly from it (common in MCUs):
  \textbf{execute-in-place} (XIP). No copy needed;
\item
  \textbf{Non-volatile}: the program is still there after power-off;
\item
  Writing is \textbf{not} a store instruction:

  \begin{itemize}
  
  \item
    First \textbf{erase} a whole sector (e.g.~a few KB) → all bits
    become 1;
  \item
    Then \textbf{program} words (bits can only go 1 → 0)
  \item
    Microseconds to milliseconds, through a special command sequence
  \end{itemize}
\item
  Result: at run time, NOR flash is effectively \textbf{read-only}.
\end{itemize}

\par\smallskip
\textbf{\textsf{NAND flash and disks}}

\begin{itemize}

\item
  NAND flash (SD cards, SSDs) and hard disks are \textbf{not} on the
  memory bus: they are reached through a controller, one \textbf{block}
  (page) at a time, e.g.~512 B -- 4 KB;
\item
  A block has no CPU address → the CPU \textbf{cannot fetch
  instructions} from it, nor read a single variable;
\item
  Contents must first be \textbf{copied into RAM}, then used from there;
\end{itemize}

This is what a desktop PC does: the program sits on the SSD, and the
OS \textbf{loader} copies it into RAM before running it. A
microcontroller has no loader: its program already sits in NOR flash, at
the address it runs from.

\subsubsection{Effect on program sections}

The different memory types we have available play a role in how we define our \textit{program headers} (or \textit{segment} headers) for embedded programs.

\begin{table}[ht]
    \centering
    \begin{tabular}{lccccc}
        \hline
        Technology & On the bus & Executable & Writable at runtime & Persistent \\
        \hline
        SRAM         & Yes & Yes (XIP) & Yes, bytes & No  \\
        NOR flash    & Yes & Yes (XIP) & No (erase/program) & Yes \\
        NAND / disk  & No (blocks) & No & Via a driver & Yes \\
        \hline
    \end{tabular}
\end{table}

The CPU sees one 4 GB address space. On every Cortex-M the \textbf{regions} are fixed by the architecture; each chip decides what it populates.

\begin{table}[ht]
    \centering
	\begin{tabular}{llp{5cm}}
        \hline
        Address range & Region & On \texttt{mps2-an385} \\
        \hline
        \texttt{0x0000\_0000}--\texttt{0x1FFF\_FFFF}
            & Code
            & 4 MB at \texttt{0x0}: the program ("FLASH") \\

        \texttt{0x2000\_0000}--\texttt{0x3FFF\_FFFF}
            & SRAM
            & 4 MB at \texttt{0x2000\_0000}: "RAM" \\

        \texttt{0x4000\_0000}--\texttt{0x5FFF\_FFFF}
            & Peripherals
            & UART, timers\ldots \\

        \texttt{0x6000\_0000}--\texttt{0xDFFF\_FFFF}
            & External RAM / devices
            & Not used here \\

        \texttt{0xE000\_0000}--\texttt{0xE00F\_FFFF}
            & CPU's private registers
            & Control of the CPU itself (at \texttt{0xE000\_E000}) \\
        \hline
    \end{tabular}
    \caption{Cortex-M memory map and the regions populated by \texttt{mps2-an385}.}
    \label{tab:cortex-m-memory-map}
\end{table}

\subsubsection{Linker script}

This leaves us with the following linker script, which we can find very useful to analyze:
\begin{lstlisting}[language=C, style=codestyle]	
/* Linker script for mps2-an385 machine (ARM Cortex-M3, MPS2 platform). */

/* Physical memory regions of the board. */
MEMORY
{
    /* code, rodata (XIP) */
    FLASH (rx)  : ORIGIN = 0x00000000, LENGTH = 4M

    /* data, bss, stack */
    RAM   (rwx) : ORIGIN = 0x20000000, LENGTH = 4M
}

/* Entry point, also entry 1 of vector table (reset handler). */
ENTRY(Reset_Handler)

/* initial SP = top of RAM; the stack grows downward from here */
_estack = ORIGIN(RAM) + LENGTH(RAM);  

/* 4 KB reserved for the stack*/
_Min_Stack_Size = 0x1000;             

/* combine input (ELF) sections into output sections */
SECTIONS
{
    /* Vector table first, at address 0, so the CPU finds the initial SP
     * and reset vector where the architecture requires them. */
    .isr_vector :
    {
        . = ALIGN(4);            /* vector table must be word-aligned */
        KEEP(*(.isr_vector))     /* KEEP: never garbage-collect it  */
        . = ALIGN(4);            /* keep alignment */
    } > FLASH                    /* place in FLASH */

    /* Executable code and read-only data. */
    .text :
    {
        . = ALIGN(4);           /* word-align the start of this section */
        *(.text*)               /* all code from every input object */
        *(.rodata*)             /* string literals, const tables */
        . = ALIGN(4);           /* word-align the end so (_etext) */
        _etext = .;             /* end-of-text marker */
    } > FLASH                   /* place in FLASH */
    
    /* FLASH address where .data's initial image is stored: the technique
    allows for reserving space in RAM while storing in flash. */
    _sidata = LOADADDR(.data);  

    /* Initialized variables: they live (run) in RAM but their initial
     * contents are stored in FLASH and copied over by Reset_Handler. */
    .data :
    {
        . = ALIGN(4);           /* word-align the start of .data */
        _sdata = .;             /* start of .data in RAM */
        *(.data*)               /* every initialized-variable input section */
        . = ALIGN(4);           /* word-align the end */
        _edata = .;             /* end of .data in RAM - copy loop stops here */
    } > RAM AT > FLASH          /* place in FLASH, access in RAM */

    /* Zero-initialized variables: reserve RAM, no stored image. */
    .bss :
    {
        . = ALIGN(4);           /* word-align the start of .bss */
        _sbss = .;              /* start of .bss */
        *(.bss*)                /* every zero-initialized input section */
        *(COMMON)               /* legacy tentative-definition symbols */
        . = ALIGN(4);           /* word-align the end */
        _ebss = .;              /* end of .bss - zero loop stops here */
    } > RAM
}
\end{lstlisting}

For now we intentionally skip the \textit{heap} and other advanced features, but this should be pretty useful to understand how the memory is laid out.
Let us remember that here we are defining memory \textbf{regions}, declaring memory \textbf{sections}, and expecting the linker to generate for us memory \textbf{segments} (program headers) which place the data in our segments within more "blocks" (that is, segments) to load in different memory regions (the ones seen above).

\subsubsection{LMA vs VMA}

Let us also note the \textit{LMA/VMA} distinction:
\begin{itemize}
	\item \textbf{LMA} (\textit{Load Memory Address}) is where the initial bytes of data are located in the executable/image (the program headers which we load);
	\item \textbf{VMA} (\textit{Virtual/Run-time Memory Address}) is where the data in the section expects to live when the program executes.
\end{itemize}

The gimmick at lines 45-56 allows us to set a LMA for the data section inside the flash memory, and set its VMA to the address it will be copied to in RAM.
Of course, the task of copying the data itself is assigned to the reset handler.

We also note that the \texttt{.bss} initialize sector has a VMA (where to access uninitialized data) but no LMA (nothing to load, the reset handler just zeroes out).

\subsubsection{Vector table}

One important region of the memory (more specifically, FLASH memory) is the \textbf{vector table}.

\begin{itemize}
	\item It is a \textbf{table} of function addresses at a fixed place (address
    \texttt{0x0} at reset).
    
    \item Each \textbf{exception} (an event that makes the CPU stop and
    run a specific handler: reset, a fault, \ldots) has a numbered
    \textbf{slot}; the CPU reads slot $N$ to find where to jump,
    
    \item Slot 0 is special: it contains not an address of code, but the
    initial stack pointer (SP), which gets loaded automatically by the CPU. This means that the reset handler is at index 1 (not 0).
\end{itemize}

In \texttt{boot/startup.s} (base), section \texttt{.isr\_vector}:

\begin{lstlisting}[language={arm}]
isr_vector:
    .word _estack          /* 0: initial SP  (C: &_estack)      */
    .word Reset_Handler    /* 1: reset       (C: Reset_Handler) */
\end{lstlisting}

The stored value of word 1 is \texttt{Reset\_Handler | 1}: the low bit
marks Thumb code. Nothing ever \emph{calls} \texttt{isr\_vector}; only
the hardware reads it.

\subsection{Registers}

We now go over the ARM architecture programmer registers.
In all ARM processors, the following registers are available and accessible in any processor mode:
\begin{itemize}
	\item 13 general-purpose registers: \texttt{R0-R12}: user variables, with call arguments passed through \texttt{R0-R3};
	\item One \textit{Stack Pointer} (\texttt{SP}): address of the top of the stack;
	\item One \textit{Link Register} (\texttt{LR}): where to return after a function call;
	\item One \textit{Program Counter} (\texttt{PC}): address of the next instruction to execute;
	\item One \textit{Status Register} (\texttt{xPSR}): flags, Thumb bit, current exception.
\end{itemize}

Some things to note at this level which may not be clear:

\begin{itemize}
    \item There are no ``variables'' at this level: only these registers and
    memory.
    \item The stack grows \textbf{downwards}: a push first does
    $\mathrm{SP} \mathrel{-}= 4$, then writes.
\end{itemize}

\par\smallskip

\textbf{xPSR} packs three things into one register:

\begin{table}[ht]
    \centering
    \begin{tabular}{cl}
        \hline
        Bits & Content \\
        \hline
        31--28 & flags \textbf{N Z C V} (negative, zero, carry, overflow),
                  set by \texttt{cmp} \\
        24 & \textbf{T}: Thumb state --- must always be 1 on Cortex-M \\
        8--0 & number of the exception being handled (0 = none) \\
        \hline
    \end{tabular}
    \caption{xPSR fields}
    \label{tab:xpsr-fields}
\end{table}

The CPU is always in one of two \textbf{modes}:

\begin{itemize}
    \item \textbf{Thread mode} --- normal code: \texttt{Reset\_Handler},
    \texttt{main()} (xPSR bits 8--0 = 0).
    
    \item \textbf{Handler mode} --- running an exception handler
    (bits 8--0 = its number).
\end{itemize}

Code can also be \emph{privileged} or not; everything in this unit runs
privileged.

\subsection{ISA}

\lstset{language=arm}

ARM Cortex-M uses a \textbf{load/store} architecture: arithmetic and
logical operations operate on registers, while memory is accessed explicitly
through load and store instructions.

\begin{itemize}
    \item \texttt{ldr}: \textbf{load} data from memory into a register;
    \item \texttt{str}: \textbf{store} data from a register into memory;
    \item Arithmetic instructions do not directly operate on memory operands.
\end{itemize}

Let's have a table of example memory operations:
\begin{table}[H]
    \centering
    \begin{tabular}{ll}
        \hline
        Assembly & C equivalent \\
        \hline
        \texttt{ldr r3, [r0]} &
            \texttt{r3 = *r0;} \\

        \texttt{str r3, [r1]} &
            \texttt{*r1 = r3;} \\

        \texttt{ldr r3, [r0], \#4} &
            \texttt{r3 = *r0; r0 += 4;} \\
        & \texttt{// i.e. r3 = *src++;} \\

        \texttt{str r3, [r1], \#4} &
            \texttt{*r1 = r3; r1 += 4;} \\
        & \texttt{// i.e. *dst++ = r3;} \\

        \texttt{strb r1, [r0, \#7]} &
            \texttt{((uint8\_t *)r0)[7] = r1;} \\
        & \texttt{// one byte, at offset 7} \\
        \hline
    \end{tabular}
\end{table}

\subsection{Post- and pre-increment}

In ARM we have \textbf{pre-indexed} and \textbf{post-indexed} memory operations, that is control over the time when the base register is updated.

	      Being more specific, ARM allows load and store instructions to optionally \textit{alter} the address contained within a base register.
	      3 forms are available:
	      \begin{itemize}
		      \item \lstinline|ldr r0, [r1, #4]| is equivalent to immediate offset addressing in RISC-V, with \lstinline|#4| being the immediate operand;
		      \item \lstinline|ldr r0, [r1, #4]!|, with the additional \lstinline|!|, means performing a \textit{write-back} or \textbf{pre-indexing} operation. The calculated address, obtained by adding the offset, is \textit{written back} into the base register before the load operation. The result we get is that of a load from the base address plus the offset, with the base address reflecting the change after execution. Compared to RISC-V, this saves an instruction;
		      \item \lstinline|ldr r0, [r1], #4| means performing a \textbf{post-indexing} operation. This means increasing the value of the base register by offset, after performing the load. We can thing of this as being analogous to the \lstinline|++| syntax available in C. This also saves an instruction compared to RISC-V.
	      \end{itemize}

	      As we can see, the main point of this complication is to save a few instructions when accessing arrays, or generally where striding memory access patterns are needed.
Since a word is 4 bytes on this machine, this is useful for walking through an array or
copying a block of memory.

For example:

\begin{lstlisting}[language=arm]
copy:
    ldr r3, [r0], #4
    str r3, [r1], #4
    b   copy
\end{lstlisting}

copies one word at a time while automatically advancing both pointers.

\subsection{Loading constants and addresses}

Assembly code also needs mechanisms for loading constants, referring to
addresses, defining labels, and changing the control flow.

Let's have a table of example (pseudo and normal) instructions for constant loading.
\begin{table}[H]
    \centering
    \begin{tabular}{ll}
        \hline
        Assembly & C equivalent \\
        \hline
        \texttt{ldr r0, =\_estack} &
            \texttt{r0 = (uint32\_t)\&\_estack;} \\

        \texttt{mov sp, r0} &
            \texttt{sp = r0;} \\

        \texttt{movs r3, \#0} &
            \texttt{r3 = 0;} \\

        \texttt{copy\_data:} &
            label: a name for this address \\

        \texttt{b copy\_data} &
            \texttt{goto copy\_data;} \\

        \texttt{bl main} &
            \texttt{main();} \\
        \hline
    \end{tabular}
\end{table}

The pseudo-instruction

\begin{lstlisting}[language=arm]
ldr r0, =_estack
\end{lstlisting}

loads the value represented by \texttt{\_estack} into \texttt{r0}. This is
often an address, but the syntax can be used for arbitrary 32-bit constants.

A large constant generally cannot be encoded directly in a single Thumb
instruction. The assembler can therefore place the constant in a nearby
\textbf{literal pool} and generate a PC-relative load.

For example, GDB may show:

\begin{lstlisting}[language=arm]
ldr r0, [pc, #92]
\end{lstlisting}

The directive

\begin{lstlisting}[language=arm]
.ltorg
\end{lstlisting}

asks the assembler to emit the current literal pool at that point.

Note that a label names an address in the program:

\begin{lstlisting}[language=arm]
copy_data:
    ...
\end{lstlisting}

A branch transfers execution to that address:

\begin{lstlisting}[language=arm]
b copy_data
\end{lstlisting}

This is analogous to a C \texttt{goto}, although assembly labels are
addresses in the instruction stream rather than variables.

The instruction \texttt{bl} means \textbf{branch with link}. It branches to
the target and simultaneously saves the return address in the
\textbf{link register (LR)}. This is the normal mechanism used for a
subroutine call:

\begin{lstlisting}[language=arm]
bl main
\end{lstlisting}

A corresponding return normally uses the address stored in \texttt{lr}.

\subsubsection{Branches}

\begin{table}[ht]
    \centering
    \begin{tabular}{ll}
        \hline
        Assembly & C equivalent \\
        \hline
        \texttt{cmp r1, r2} &
            compute \texttt{r1 - r2} and update the flags \\

        \texttt{bhs done} &
            \texttt{if (r1 >= r2) goto done;} \\
        & \texttt{// unsigned comparison} \\

        \texttt{orr r1, r1, \#0x10} &
            \texttt{r1 |= 0x10;} \\

        \texttt{cpsid i} &
            disable interrupt acceptance \\

        \texttt{cpsie i} &
            enable interrupt acceptance \\
        \hline
    \end{tabular}
    \caption{Conditional branches, bit operations and interrupt control.}
    \label{tab:arm-decisions-bits-interrupts}
\end{table}

\texttt{cmp} performs a subtraction internally (\texttt{r1} - \texttt{r2}) but does not store the result. Instead, it updates the condition flags
(\texttt{N}, \texttt{Z}, \texttt{C}, and \texttt{V}) in the xPSR.

A subsequent conditional branch examines these flags.

For example, \texttt{bhs} means ``higher or same'' for an \textbf{unsigned}
comparison:

\begin{lstlisting}[language=arm]
cmp r1, r2
bhs done
\end{lstlisting}

This corresponds to:

\begin{lstlisting}[language=C]
if (r1 >= r2)
    goto done;
\end{lstlisting}

A C \texttt{while} loop therefore naturally becomes a test at the top,
followed by a backwards branch at the bottom:

\begin{lstlisting}[language=arm]
loop:
    cmp r1, r2        @ while (r1 < r2) {
    bhs done
    ...               @     body
    b   loop          @ }
done:
\end{lstlisting}

\subsubsection{Interrupt masking}

On Cortex-M, \texttt{cpsid i} disables the acceptance of maskable
interrupts, while \texttt{cpsie i} enables them again:

\begin{lstlisting}[language=arm]
cpsid i
...
cpsie i
\end{lstlisting}

These correspond conceptually to the CMSIS functions
\texttt{\_\_disable\_irq()} and \texttt{\_\_enable\_irq()}.

\subsubsection{Assembler directives}

Assembly source contains two fundamentally different kinds of lines:

\begin{enumerate}
    \item \textbf{Instructions}, which become CPU instructions.
    \item \textbf{Assembler directives}, which control how the assembler
    constructs the object file.
\end{enumerate}

Directives generally begin with a period (\texttt{.}) and do not themselves
execute on the CPU.

For example:

\begin{lstlisting}[language=arm]
.section .isr_vector,"a",%progbits
.global isr_vector

isr_vector:
    .word _estack
    .word Reset_Handler

.section .text,"ax",%progbits
.global Reset_Handler
.thumb_func

Reset_Handler:
\end{lstlisting}

Here:

\begin{itemize}
    \item \texttt{.section} selects or creates an ELF section.
    \item \texttt{.global} makes a symbol visible to the linker and other
    object files.
    \item \texttt{.word} emits a 32-bit value into the current section.
    \item \texttt{.thumb\_func} tells the assembler that the following
    function is Thumb code.
\end{itemize}

Thus, \texttt{isr\_vector} is not a function being executed by the CPU:
the assembler is simply placing two 32-bit values into the
\texttt{.isr\_vector} section.

The corresponding idea in C is:

\begin{lstlisting}[language=C]
__attribute__((section(".isr_vector")))
const uint32_t isr_vector[] = {
    (uint32_t)&_estack,
    (uint32_t)Reset_Handler
};
\end{lstlisting}

Without \texttt{.global}, an assembly symbol is local to the object file,
similar in effect to a \texttt{static} C symbol.

\subsubsection{Thumb function addresses}

Cortex-M executes Thumb instructions exclusively. Function pointers and
exception-vector entries therefore have bit 0 set to indicate Thumb state.

For a symbol such as \texttt{Reset\_Handler}, the assembler/linker can
produce an address with bit 0 set when it is declared as a Thumb function:

\begin{lstlisting}[language=arm]
.thumb_func
Reset_Handler:
\end{lstlisting}

For example, if the actual instruction address is \texttt{0x00000040}, the Thumb function pointer is represented as \texttt{0x00000041}.

The CPU uses this bit to identify the target as Thumb code; the actual
instruction address is still aligned to an even address.

\subsection{Faults}

Let's look in more detail at \textbf{faults}, so that we can understand the whole of the \textit{handlers} defined in the ISR vector table.

An \textbf{exception} is any event that interrupts the current code and
causes the CPU to execute a handler whose address is stored in the vector
table, switching to Handler mode.

\begin{table}[H]
    \centering
    \begin{minipage}[t]{0.48\textwidth}
        \vspace{0pt}
        \centering
        \begin{tabular}{c p{0.32\linewidth} p{0.38\linewidth}}
            \hline
            \# & Exception & Raised by \\
            \hline
            0 & \textit{None} & Slot holds the initial SP \\
            1 & Reset & Power-on, reset button \\
            2 & NMI & Non-maskable interrupt \\
            3 & \textbf{HardFault} & Faults that cannot be handled elsewhere \\
            4 & \textbf{MemManage} & Memory protection violations \\
            5 & \textbf{BusFault} & Bus access errors \\
            6 & \textbf{UsageFault} & Invalid instruction use, division by zero, \ldots \\
            \hline
        \end{tabular}
    \end{minipage}
    \hfill
    \begin{minipage}[t]{0.48\textwidth}
        \vspace{0pt}
        \centering
        \begin{tabular}{c p{0.32\linewidth} p{0.38\linewidth}}
            \hline
            \# & Exception & Raised by \\
            \hline
            7--10 & \textit{Reserved} & Slot must hold \texttt{0} \\
            11 & SVCall & The \texttt{svc} instruction \\
            12 & DebugMonitor & Debug events (unused here) \\
            13 & \textit{Reserved} & Slot must hold \texttt{0} \\
            14 & PendSV & A software request \\
            15 & SysTick & Built-in timer \\
            16+ & IRQs & Device interrupts \\
            \hline
        \end{tabular}
    \end{minipage}
\end{table}

At reset, only HardFault is active among the configurable fault handlers:
faults that cannot be handled by their corresponding handlers escalate to
HardFault. This corresponds to ``escalate 3'' behavior observed during debugging.

\subsubsection{Handlers}

The initial vector table contains only two words, but the Cortex-M processor
expects one word per exception vector, for a total of 16 entries in the
initial table. Slots 7--10 and 13 are reserved and must contain zero.
Every other slot must contain a valid handler address, even if no specific
handler has been implemented yet.

A default handler can be defined as follows:

\begin{lstlisting}[language={[x86masm]Assembler}]
    .word HardFault_Handler        @ slot 3, and so on up to 15

    .weak      HardFault_Handler
    .thumb_set HardFault_Handler, Default_Handler

Default_Handler:
    b Default_Handler              @ for (;;) {}
\end{lstlisting}

The equivalent declaration in C is:

\begin{lstlisting}[language=C]
void HardFault_Handler(void)
    __attribute__((weak, alias("Default_Handler")));
\end{lstlisting}

A \textbf{weak symbol} provides a default implementation that can be
overridden by a normal definition in another object file. If
\texttt{main.c} defines its own \texttt{HardFault\_Handler}, the linker
uses that implementation instead of the default.

The completed vector table occupies 16 words:

\[
    16 \times 4 = 64 \text{ bytes} = \texttt{0x40}.
\]

Consequently, if the table starts at address \texttt{0x0}, the
\texttt{Reset\_Handler} entry point moves from \texttt{0x8} to
\texttt{0x40}. Since Cortex-M executes Thumb instructions, its vector
entry contains \texttt{0x41}, with bit 0 set to indicate Thumb state.

\subsection{CPU registers}

One interesting thing about ARM is that some CPU settings live in \textit{registers} which are stored at certain \textbf{memory addresses}.

For instance, the \textbf{SCB} (\textit{System Control Block}) is a set of 32-bit registers inside the CPU, reached like memory in the system area at \texttt{0xE000E000} – \texttt{0xE000EFFF}.

Let's recall the default processor registers:
\begin{table}[H]
    \centering
    \begin{tabular}{ll}
        \hline
        Register & Purpose \\
        \hline
        \texttt{r0}--\texttt{r12} & General-purpose registers \\
        \texttt{r13} (\texttt{SP}) & Stack pointer \\
        \texttt{r14} (\texttt{LR}) & Link register; return address \\
        \texttt{r15} (\texttt{PC}) & Program counter \\
        \texttt{xPSR} & Flags, execution state and exception number \\
        \texttt{MSP} & Main stack pointer \\
        \texttt{PSP} & Process stack pointer \\
        \texttt{PRIMASK} & Masks configurable interrupts \\
        \texttt{BASEPRI} & Masks interrupts below a priority threshold \\
        \texttt{FAULTMASK} & Masks exceptions except NMI \\
        \texttt{CONTROL} & Selects stack and privilege state in Thread mode \\
        \hline
    \end{tabular}
\end{table}

The exact register set depends on the Cortex-M implementation. In
particular, \texttt{BASEPRI} is not implemented on every Cortex-M core.

The \texttt{SP} register refers to the currently selected stack pointer,
either \texttt{MSP} or \texttt{PSP}. The processor uses the main stack
pointer in Handler mode.

These peripherals expose registers through fixed memory addresses. They
are accessed using ordinary load and store instructions.

Let's look at the main peripherals we have:
\begin{table}[H]
    \centering
    \begin{tabular}{lll}
        \hline
        Block & Base address & Purpose \\
        \hline
        SCnSCB & \texttt{0xE000E000} & System control and feature registers \\
        SysTick & \texttt{0xE000E010} & 24-bit system timer \\
        NVIC & \texttt{0xE000E100} & Interrupt enable, pending state and priorities \\
        SCB & \texttt{0xE000ED00} & Exceptions, faults and processor configuration \\
        MPU & \texttt{0xE000ED90} & Memory protection (optional) \\
        DWT & \texttt{0xE0001000} & Watchpoints, cycle counter and profiling \\
        FPB & \texttt{0xE0002000} & Flash patch and breakpoint unit \\
        ITM & \texttt{0xE0000000} & Instrumentation and debug output \\
        CoreDebug & \texttt{0xE000EDF0} & Debug control and status \\
        TPIU & \texttt{0xE0040000} & Trace output interface \\
        \hline
    \end{tabular}
\end{table}

These addresses correspond to the Cortex-M3 memory map. Optional blocks
may be absent or incompletely implemented in a particular processor or
emulator.

The \textit{System Control Block} (\textbf{SCB}) provides registers for processor
configuration, exception handling and fault diagnosis.

Let's look at some important elements of the SCB:
\begin{table}[H]
    \centering
    \begin{tabular}{lll}
        \hline
        Register & Offset & Purpose \\
        \hline
        \texttt{CPUID} & \texttt{0x00} & Processor identification \\
        \texttt{ICSR} & \texttt{0x04} & Exception and interrupt state \\
        \texttt{VTOR} & \texttt{0x08} & Vector table base address \\
        \texttt{AIRCR} & \texttt{0x0C} & System reset and priority configuration \\
        \texttt{SCR} & \texttt{0x10} & Sleep behavior \\
        \texttt{CCR} & \texttt{0x14} & Processor configuration \\
        \texttt{SHPR1--3} & \texttt{0x18--0x23} & System exception priorities \\
        \texttt{SHCSR} & \texttt{0x24} & System handler enable and status \\
        \texttt{CFSR} & \texttt{0x28} & Configurable fault status \\
        \texttt{HFSR} & \texttt{0x2C} & HardFault status \\
        \texttt{DFSR} & \texttt{0x30} & Debug fault status \\
        \texttt{MMFAR} & \texttt{0x34} & Memory-management fault address \\
        \texttt{BFAR} & \texttt{0x38} & Bus fault address \\
        \hline
    \end{tabular}
\end{table}

\subsection{Real chips}

QEMU's CPU is ready as soon as it leaves reset. A real MCU's startup
code usually does more, often in a vendor function \texttt{SystemInit()}:

\begin{itemize}
    \item \textbf{Clock / PLL}: the chip starts on a slow internal
    oscillator; the PLL multiplies the clock frequency up to the desired
    operating frequency.

    \item \textbf{Flash wait states}: Flash memory is slower than a fast
    CPU. The Flash controller must be configured to insert additional
    wait cycles \emph{before} increasing the CPU clock frequency.

    \item \textbf{Watchdog}: a timer that resets the chip unless software
    periodically ``feeds'' it. On some chips the watchdog starts running
    immediately after reset, so the startup code must feed or disable it.

    \item \textbf{FPU (Cortex-M4F)}: the floating-point unit is disabled
    at reset. It can be enabled through the \texttt{CPACR} register at
    \texttt{0xE000ED88}; otherwise executing an FPU instruction can cause
    a UsageFault.

    \item \textbf{C++}: before entering \texttt{main()}, the startup code
    must call the static constructors stored in \texttt{.init\_array}.
\end{itemize}

\end{document}
